\documentclass[12pt]{article}
\usepackage[utf8]{inputenc}
\usepackage[T1]{fontenc}
\usepackage{hyperref}
\title{Wumbvu Language Resources for K-12 Curriculum Development}
\author{Compiled by Claude Sonnet 5.5 (claude-sonnet-5-5) for Andr\'as Kornai}
\date{}
\begin{document}
\maketitle

\section{Introduction}
Wumbvu (also Wumvu or Wumbu; ISO 639-3 wum) is a Bantu language of Gabon and the Republic of Congo. Speaker estimates vary by source: Wikipedia cites 18,000 in Gabon, Joshua Project lists 45,000 in Gabon and 13,000 in Congo, and ASJP gives 28,900. Little published material turned up in the searches.

\section{Primary Educational Resources}
No primer, textbook or school programme for Wumbvu turned up in the searches.

\section{Linguistic and Reference Materials}
Reference pages include Wikipedia (a stub), \url{https://en.wikipedia.org/wiki/Wumbvu_language}, and the Glottolog entry, \url{https://glottolog.org/resource/languoid/id/wumb1242.bigmap.html}. ASJP has a wordlist entry at \url{https://asjp.clld.org/languages/WUMBVU_3}. The searches turned up no dictionary or grammar.

\section{Religious and Community Resources}
Joshua Project's language page is \url{https://www.joshuaproject.net/index.php/languages/wum}. It reports no complete New Testament and no audio Bible. Its Gabon people-group profile is at \url{https://www.joshuaproject.net/index.php/people_groups/15421/GB}. The Jesus Film is reported to exist in Wumvu.

\section{Digital and Audio Resources}
The Global Recordings Network language page is \url{https://globalrecordings.net/en/language/wum}. It appeared to list no recordings of its own, and it files the language under a different name.

\section{Curriculum Development Context}
The language is not known to be taught in schools, and written material is nearly absent. Curriculum work would start from documentation and orthography development, in cooperation with the Gabonese and Congolese communities.
\end{document}